\section{オペレーション手順}

\subsection{立ち上げシーケンス}
MCP保護のため、必ず以下の順序を遵守すること。
\begin{enumerate}
    \item **真空確認:** MCPチャンバーの圧力が $1.0 \times 10^{-4}\,\mathrm{Pa}$ 以下であることを確認。
    \item **冷却確認:** TMPおよびカメラのファンが回っていることを確認。
    \item **制御系ON:** パルスジェネレータ、カメラコントローラの電源ON。
    \item **HV印加:** \begin{itemize}
        \item まずMCP-in用電源（負電圧）を印加。
        \item 次にPhosphor用電源（正電圧）を印加。
        \item \textbf{警告:} 真空度が悪い状態でHVを印加すると、MCP内で放電が起き一撃で破損する。
    \end{itemize}
\end{enumerate}

\subsection{実験中の確認事項}
トリガータイミングがずれると、ノイズだらけの放電開始時（Breakdown）を撮影してしまい、高電圧回路に過剰な負荷がかかる。オシロスコープで常にゲートタイミングを監視すること。